
% resumo em português
\begin{resumo}[Resumo] 
\noindent Texto de resumo em português, com 150 a 500 palavras. O resumo deve ter apenas 1 parágrafo, sem citações e sem recuo de texto na primeira linha. Também não pode haver quebra de linha. O resumo em português é necessário mesmo que o seu texto principal esteja em inglês. Infelizmente, voltamos aos trechos que você não pode simplesmente deixar de escrever, uma vez que são obrigatórios para o trabalho. Contudo, o resumo é um trecho bem rápido e fácil de escrever, não deve dar dor de cabeça. Geralmente é aqui que o leitor deixa a preguiça vencer ou continua perseverando para ler o seu trabalho. Mas não é o espaço para convencer o leitor a permanecer, ainda. Apenas para apresentar sobre o que se trata o trabalho de forma breve, sucinta. Se atente para seguir a formatação correta, de apenas 1 parágrafo, sem recuo de texto na primeira linha, nem quebras de linha.
% \noindent %- o resumo deve ter apenas 1 parágrafo e sem recuo de texto na primeira linha, essa tag remove o recuo. Não pode haver quebra de linha.

 \vspace{\onelineskip}
    
 \noindent
 \textbf{Palavras-chave}: Palavras-chave; separadas; por; ponto; e; virgula.
\end{resumo}



% resumo em inglês
\begin{resumo}[Abstract]
\begin{otherlanguage*}{english}

 %\noindent
\noindent Abstract text in English, with 150 to 500 words. The summary must be in a single paragraph, with no indentation at the beginning of the line. There should also be no line breaks. The summary in portuguese is required even if your main text is in english. Unfortunately, we return to the sections that you simply cannot skip, as they are mandatory for the paper. However, the abstract is a very quick and easy section to write and should not cause any headaches. It is usually here that the reader either gives in to laziness or keeps going and perseveres in reading your work. But this is not the space to persuade the reader to continue just yet. It is only meant to briefly and concisely present what the work is about. Be sure to follow the correct formatting: a single paragraph, with no indentation at the beginning and no line breaks.

   \vspace{\onelineskip} 
 
   \noindent 
      \textbf{Keywords}: Keywords; separated; by; semicolon; Text; Text.

  \end{otherlanguage*}
 \end{resumo}
